% Part: many-valued-logic
% Chapter: infinite-valued-logics
% Section: introduction

\documentclass[../../../include/open-logic-section]{subfiles}

\begin{document}

\olfileid{mvl}{inf}{int}

\olsection{Pambuka}

Cacah nilai kabeneran ing sawijining matriks ora kudu winates.
Pilihan kang cetha kanggo himpunan nilai kabeneran tanpa wates
yaiku himpunan wilangan rasional antarane $0$ lan~$1$,
$V_\infty = [0,1] \cap \Rat$, yaiku,
\begin{align*}
    V_\infty & = \Setabs{\frac{n}{m}}{n,m \in \Nat, 0<m, n\le m}.
\intertext{Nalika nimbang himpunan nilai kabeneran tanpa wates iki,
asring migunani uga nimbang subhimpunan kanggo indeks paling ora loro:}
V_m & = \Setabs{\frac{n}{m-1}}{n \in \Nat \text{ lan } n<m}
\intertext{Umpamane, $V_5$ iku himpunan kang dumadi saka $5$ nilai kabeneran
kanthi sela kang padha,}
V_5 & = \{0, \frac{1}{4}, \frac{1}{2}, \frac{3}{4}, 1\}.
\end{align*}
Ing logika kang adhedhasar himpunan nilai kabeneran iki, lumrahe
mung $1$ kang dadi nilai pinilih, yaiku $V^+ = \{1\}$. Kanthi
tembung liya, $1$ dadi nilai bener (mutlak), $0$ dadi nilai salah
mutlak, nanging !!{formula}s bisa nduweni nilai antarané
ing~$V$.

Kita uga bisa nimbang himpunan $V_{[0,1]} = [0,1]$ kang ngemot
kabeh wilangan \emph{real} antarane $0$ lan~$1$, utawa
subhimpunan tanpa wates liyane saka $[0,1]$. Logika kang
nggunakake himpunan nilai kabeneran iki kerep diarani
\emph{fuzzy}.

\end{document}
